\documentclass[a4paper,12pt]{article}
\usepackage[margin=1.5cm]{geometry}
\usepackage{amsmath, amssymb}
\usepackage{array}
\usepackage{multicol}
\usepackage{graphicx}

\begin{document}
\subsection*{Tema 3: Fracciones}
    \begin{enumerate}
        \item Representa las siguientes fracciones: 
            \begin{enumerate}
                \item $\dfrac{3}{8}$
                \item $\dfrac{3}{4}$
                \item $\dfrac{5}{4}$
                \item $\dfrac{7}{3}$
            \end{enumerate}        
        \item Calcula:
            \begin{enumerate}
                \item $\dfrac{2}{3}$ de 36
                \item $\dfrac{8}{9}$ de 72
                \item $\dfrac{7}{6}$ de 56
                \item $\dfrac{3}{4}$ de 54 
            \end{enumerate}
        \item Comprueba si las siguientes fracciones son equivalentes.
            \begin{enumerate}
                \item $\dfrac{2}{3} y \dfrac{5}{7}$
                \item $\dfrac{4}{6} y \dfrac{6}{9}$
                \item $\dfrac{1}{3} y \dfrac{2}{9}$
                \item $\dfrac{2}{15} y \dfrac{5}{8}$
            \end{enumerate}
        \item Ordena de menor a mayor
            \begin{enumerate}
                \item $\dfrac{7}{15} , \dfrac{9}{20}$
                \item $\dfrac{8}{3} , \dfrac{3}{2}, \dfrac{9}{5}$
                \item $\dfrac{3}{10}, \dfrac{-8}{15}, \dfrac{12}{25}$
            \end{enumerate}
        \item Realiza las siguientes sumas
            \begin{enumerate}
                \item $\dfrac{3}{10}+\dfrac{2}{10}$
                \item $\dfrac{5}{8} + \dfrac{7}{8}$
                \item $\dfrac{3}{10}-\dfrac{2}{5}+\dfrac{3}{2}$
                \item $\dfrac{4}{5}+\dfrac{2}{7}-1$
            \end{enumerate}
        \item Realiza las siguientes multiplicaciones y divisiones, reduciendo al máximo su resultado
            \begin{enumerate}
                \item $\dfrac{2}{7} \cdot \dfrac{21}{16}$
                \item $\dfrac{12}{7} \cdot   2$
                \item $\dfrac{-2}{3} : \dfrac{4}{5}$
                \item $\dfrac{-12}{30}: \dfrac{14}{35}$
            \end{enumerate}
        \item Realiza las siguientes operaciones combinadas
            \begin{enumerate}
                \item $\dfrac{2}{7} \cdot \dfrac{21}{5} - \dfrac{3}{5} \cdot   \dfrac{15}{6}$
                \item $\dfrac{3}{5}+ (\dfrac{7}{15}-1)$
                \item $\dfrac{5}{12}: (\dfrac{3}{4}- \dfrac{2}{3})$
                \item $(1- \dfrac{5}{12}): (\dfrac{1}{4} - \dfrac{2}{3})$
            \end{enumerate}
        \item ¿Cuantas botellas de $\dfrac{3}{4}$ de litro se pueden llenar con una garrafa de 30 litros?
        \item Alberto ha fallado 2 penaltis de 5 y Carlos ha marcado 4 de 7. ¿Quién tira mejor los penaltis?
        \item El suelo de un almacén tiene 1200 m² de superficie. Luis pinta un día una cuarta parte, y otro día un tercio de dicha superficie. Su compañero Juan pinta el resto.
        \begin{enumerate}
            \item ¿Qué fracción del almacén pinta cada uno de ellos?
            \item Si pagan dos euros el metro cuadrado, ¿Cuánto gana cada uno?
        \end{enumerate}
        \item Invité a mis primos los zampabollos a mi fiesta de cumpleaños. Tienen fama de comerse todo lo que les pasa por delante. Tras soplar las velas fui un momento al baño y, al volver, se habían comido 5/8 de mi tarta de cumpleaños. Puse lo que quedaba de tarta en la báscula y esta marcó 600 gramos. ¿Cuánto pesaba la tarta inicialmente? Representa gráficamente la situación.
    \end{enumerate}
\subsection*{Unidad 4: Decimales}
    \begin{enumerate}
        \item Escribe como se lee los siguientes números decimales:
        \begin{enumerate}
            \item 0.015
            \item 1.005
            \item 12.7
            \item -10.55
        \end{enumerate}
        \item Ordena de menor a mayor los siguientes números
            \begin{enumerate}
                \item 2.7 , 2.691, 2.71, 2.689 y 3
                \item -2.07, -2.6091, -2.071, -2.6089 y -3.09
            \end{enumerate}
        \item Completa la tabla
            \begin{tabular}{| c | c | c |}
                \hline
                  & Redondea a las décimas & Redondea a las centésimas \\
                \hline
                15.789&&\\ 
                \hline
                20.234&&\\ 
                \hline
                12.794&&\\ 
                \hline
                4.8412&&\\ 
                \hline
            \end{tabular}
        \item Efectúa las siguientes sumas y restas
            \begin{enumerate}
                \item 21.756+23.89=
                \item 19.4 + 28.78 - 30.791=
            \end{enumerate}
        \item Realiza las siguientes multiplicaciones con números decimales:
            \begin{enumerate}
                \item $5.2 \cdot 7.02 = $
                \item $10.25 \cdot 8.7 = $
            \end{enumerate}
        \item Realiza las siguientes divisiones con números decimales, calcula dos cifras decimales e indica el cociente y el resto. 
            \begin{enumerate}
                \item 15.87 : 5 = 
                \item 0.23 : 1.2 =
            \end{enumerate}
        \item Realiza las siguientes operaciones con números decimales.
            \begin{enumerate}
                \item $0.00075  \cdot 1000= $
                \item $-1.5 \cdot (-1000)=$
                \item $18.75 : (-10)=$
                \item $-250.15 : (-100)=$
            \end{enumerate}
        \item Realiza las siguientes operaciones combinadas.
            \begin{enumerate}
                \item $(2-4.5-1.4) \cdot (-2.15-7.4) = $
                \item $(1:4+1.7 \cdot2 -5) \cdot(100)=$
            \end{enumerate}
        \item Laura va a una frutería y compra 2 kg de naranjas que le cuesta 1.75 € el Kg y 1.8 Kg de tomates a 2.95 € el Kg. Al final paga con un billete de 20€, ¿cuánto dinero le ha sobrado a Laura en la frutería?
        \item Juan paga con 20€ un libro (12.45€) y estuche (3.80€). ¿Cuánto dinero le sobra?
    \end{enumerate}
\subsection*{Tema 5: Álgebra}
    \begin{enumerate}
        \item Escribe en lenguaje algebraico las siguientes expresiones y calcula el resultado:
            \begin{enumerate}
                \item La mitad de un número más cuatro unidades 
                \item El cuadrado de un número menos dos.
            \end{enumerate}
        \item  Dados los polinomios 
            \[P(x) = 2x^5+3x-2  \qquad \qquad Q(x) = 3x^5+x^2-7x+1\]
            Calcula: \\
                (a) P(x)+Q(x) \\
                (b) Q(x)-P(x)
            $\vspace{5cm}$ 
        \item  Resuelve las siguientes ecuaciones:
            \begin{enumerate}
                \item $8x-5x=x+8  $
                \item $2x-8=1-3(x-2)$ 
                \item $11+2x=6x-3+3x$
                \item $5x=4-3x+5-x$
                \item $5 x - ( 4 - 2 x ) = 2 - 2 x$
                \item $1 - 6 x = 4 x - ( 3 - 2 x )$
                \item $6 ( x - 2 ) - x = 5 ( x - 1 )$
                \item $3 ( 3 x - 2 ) - 7 x = 6 ( 2 x - 1 ) x$
            \end{enumerate}
            
       
        \item  Completa la tabla para cada uno de los monomios siguientes.
        
            \begin{tabular}{| c | c | c | c |}
                \hline
                Monomio & Coeficiente & Parte literal & Grado \\
                \hline
                $2a$&&&\\ 
                \hline
                $x^2$&&&\\ 
                \hline
                $-3ab$&&&\\ 
                \hline
                $\dfrac{1}{2}xy^3$&&&\\ 
                \hline
                $-7y^3$&&&\\ 
                \hline
            \end{tabular}
        
        \item  Determina los miembros, los términos ,incógnitas y el grado de la ecuación.
        
        \begin{tabular}{| c | c | c | c | c | c |}
                \hline
                Ecuación & 1º miembro & 2º miembro & Términos & Incógnita & Grado \\
                \hline
                $7x-3 = 3x^2 +9$ &&& $\qquad \qquad \qquad \qquad $ &&\\ \hline
                $3x^^2+7x-1 = 2-3x^5+3x$ &&& $\qquad \qquad \qquad \qquad $ &&\\ 
                \hline
            \end{tabular}
           
        \item Efectúa las siguientes sumas y restas de polinomios. 
            \begin{enumerate}
                \item $3 x^2 + 2 x^2 =$
                \item $5x^4+6x^4+x^4+2x^2=$
                \item $6x^3+7x^3-10x^3=$
            \end{enumerate}

        \item Halla el valor numérico de los siguientes polinomios
        \begin{enumerate}
            \item $P(x) = 4x^5-6x^4+3x^2+9x-8$ Para x =0
            \item $Q(x) = -x^3+2x-3 $ Para x=-1
            \item $B(x) = x^2-x+2$ Para x=-3
        \end{enumerate}
        \item  Luis tiene dos años más que su hermana y entre los dos suman 18 años. ¿Cuál es la edad de cada hermano? 
        \item  La suma de tres números consecutivos es 30. Hállalos. 
        \item Si al doble de un número se le resta su mitad, resulta 54. Hallar ese número.
        \item Un kilo de cerezas cuesta dos euros más que uno de peras. Natalia ha pagado 8 € por 3 kilos de peras y uno de cerezas. ¿A cómo están las unas y las otras?
    \end{enumerate}
\end{document}